\chapter*{About This Monograph}
\addcontentsline{toc}{chapter}{About This Monograph}
\markboth{About This Monograph}{About This Monograph}

Turning a biological product from a living culture into a finished, specified
material is a chain of very different operations: a \emph{fermentation} that grows
the molecule, a sequence of \emph{separations} that purifies it, and the physical
measurements---viscosity, texture, pressure drop---that characterise it along the
way. Each is expensive to run, each obeys its own mechanism, and each is a place
where a well-chosen \emph{experiment} saves weeks and a badly chosen one wastes
them. This is exactly the setting in which \emph{Design of Experiments} (DoE)
earns its keep---and exactly the setting this monograph builds a virtual
laboratory for. Among the process operations, \emph{preparative chromatography} is
the workhorse and, mathematically, the richest: a chromatographic column is a
distributed-parameter system in which convection, dispersion, and nonlinear,
multi-component adsorption interact, and whose behaviour swings by orders of
magnitude as a handful of mobile-phase variables are turned. It is treated here in
depth, but it is one system among several.

This monograph documents the scientific content of the \code{VlabDOE}
(Virtual Lab DOE) software package: a \emph{virtual laboratory} for bioprocess
development and for the simple benchtop systems on which experimental design is
best learned. The package does two things. First, it provides mechanistic,
first-principles models---of preparative chromatography, of
ultrafiltration/diafiltration, of a living milk fermentation, and of a handful of
benchtop chemical-engineering experiments---that act as a synthetic ``ground
truth'': an oracle that, when probed with a set of process conditions, returns the
chromatogram, pH curve, or instrument reading that a real experiment would.
Second, it provides the statistical and machine-learning machinery of
\emph{Design of Experiments} (DoE) that one uses to interrogate such an oracle
efficiently: factorial designs, space-filling designs, latent-variable analysis,
covering-array screening, and adaptive Bayesian optimisation.

The document is organised into three parts. \textbf{Part~I (Foundations)} is for the
reader who arrives strong in one of this book's two parent disciplines but new to
the other. Because the monograph fuses statistics with process and separation
science, it is easy to assume of either field a fluency the reader may not yet
have; these two chapters supply it, each built from the ground up but curated to
exactly the ideas the later chapters use.

\begin{itemize}
  \item \textbf{Chapter~\ref{ch:stats}} introduces the statistical vocabulary---
  distributions, estimation, least squares and the analysis of variance,
  cross-validation, the bootstrap, the Bayesian update, and the crucial
  distinction between \emph{variability} and \emph{uncertainty}---that the design
  methods rest on.

  \item \textbf{Chapter~\ref{ch:sep}} introduces the separation science---unit
  operations, the molecules and their impurities, the adsorption isotherm, the
  column and its plates and resolution, mass transfer, the modulator idea that
  unifies the chromatography modes, and the basics of microbial fermentation.
\end{itemize}

\textbf{Part~II (The Virtual Laboratory)} is the technical core.

\begin{itemize}
  \item \textbf{Chapter~\ref{ch:models}} develops the chromatography models. We
  begin from the governing conservation laws, reduce them to the equations the
  package actually integrates, and then treat each adsorption mode---ion
  exchange, hydrophobic interaction, and reversed phase---in turn. For every
  model we give the full mathematical formulation, a careful description of each
  parameter and its units, the real-world separation problems the model
  represents, the assumptions and limitations that bound its validity, and the
  experiments and data one performs to estimate its parameters.

  \item \textbf{Chapter~\ref{ch:design}} develops the experimental-design
  methods. Each method---full factorial designs with response-surface analysis,
  Latin hypercube sampling, principal-component and partial-least-squares
  analysis, and Gaussian-process Bayesian optimisation---is presented with its
  theoretical underpinning, its mathematical formulation, and an honest
  accounting of where it is strong and where it fails.

  \item \textbf{Chapter~\ref{ch:ferment}} turns to a different kind of unit
  operation---a living \emph{milk fermentation}, the basis of yogurt and cultured
  dairy. Here the model is mechanistic but irreducibly stochastic, the only
  observable is an indirect pH curve, and the central design variable is
  \emph{combinatorial}: which of many candidate strains to blend. The chapter
  develops the fermentation model and its three layers of randomness, the
  real-world experiment, the \emph{covering-array} and related designs for
  screening large strain libraries, and the tree-ensemble methods used to analyse
  the resulting data.
\end{itemize}

\textbf{Part~III (A Benchtop Sandbox)} steps deliberately down in complexity, to
systems whose mechanism is one or two equations rather than a system of coupled
PDEs.

\begin{itemize}
  \item \textbf{Chapter~\ref{ch:benchtop}} assembles four \emph{benchtop}
  systems---pressure drop in a pipe, a falling-ball viscometer, back-extrusion of
  a set yogurt, and the hydrolysis of methyl acetate---each an experiment a
  student could run in an afternoon. Between them they exhibit the response
  \emph{shapes}---power law, calibration, threshold, and reaction kinetics---that
  most often decide which design to reach for, making them a sandbox in which the
  behaviour of an experimental design is laid bare with nothing to hide behind.
\end{itemize}

Throughout, the mathematics is tied back to the code. Where a symbol in an
equation corresponds to an argument of a function or a field of a data
structure, we say so, in \code{typewriter} font. The intent is that a reader can
move from the page to the source and back without losing the thread.

A word on philosophy. The models in this package are deliberately the
\emph{simplest models that retain the right physics}. The linearised
Steric-Mass-Action isotherm, the lumped transport-dispersive band-broadening
term, the film model of ultrafiltration flux: each is a reduction of a more
complete theory, chosen because it captures the dependence of the response on
the critical process parameters faithfully while remaining cheap enough to
evaluate thousands of times inside a design loop. We have tried to be explicit
about every such reduction, because knowing the boundary of a model's validity
is the difference between a tool and a trap.

\vspace{1em}
{\itshape\noindent
The figures in this monograph are generated directly from the package source by
\code{doc/scripts/make\_figures.py} (and, for the foundations chapters,
\code{make\_foundations\_figures.py}); the chromatograms are real numerical
solutions of the models described in the text, and even the introductory figures
are computed from the small teaching module
\code{vlab\_doe.foundations}---not schematics.\par}
